\documentclass[10pt,a4paper]{amsart}
\usepackage[margin=19mm]{geometry}
\usepackage{amsmath,amssymb,mathtools}
\usepackage[scr=rsfs,cal=euler]{mathalfa}
\usepackage{xfrac}
\usepackage[unicode,hidelinks,bookmarks=false]{hyperref}
\newcommand{\ie}{i.e.\ }
\newcommand{\cf}{cf.\ }
\newcommand{\resp}{respectively\ }
\newcommand{\Subsection}{\subsection}
\providecommand{\nobreakdash}{\nobreak\hbox{-}}
\setlength{\emergencystretch}{2em}
\multlinegap0pt
\usepackage[shortlabels]{enumitem}
\newtheorem{proposition}{Proposition}
\newtheorem{theorem}{Theorem}
\DeclareMathOperator{\K}{Gr}
\DeclareMathOperator{\bi}{bi}
\DeclareMathOperator{\irr}{Irr}
\newcommand{\lf}{\mathfrak{L}}
\DeclareMathOperator{\lie}{Lie}
\title{Tensor products of Leibniz bimodules and Grothendieck rings}
\author{J\"org Feldvoss, Friedrich Wagemann}
\date{Source: arXiv:2411.01044v4 (CC BY 4.0)}
\begin{document}
\maketitle

\noindent\emph{Source and licence.} Tensor products of Leibniz bimodules and Grothendieck rings. Version arXiv:2411.01044v4 (CC BY 4.0), distributed by arXiv under the Creative Commons Attribution 4.0 International licence, http://creativecommons.org/licenses/by/4.0/, as recorded in the arXiv licence metadata for this exact version and shown on the abstract page https://arxiv.org/abs/2411.01044v4. Full licence notice: \texttt{licenses/arxiv-2411.01044-CC-BY-4.0.txt}.\par\smallskip

\noindent\textit{Editorial scope.} Counted unit: the article's theorem alt (source0.tex lines 1668--1670). The article's complete original proof of it is delivered with it (source0.tex lines 1672--1759, one \texttt{proof} environment of 88 lines). No mathematics is written, repaired or re-derived: the statement and the proof are the article's own lines, in the article's own notation, and no retained line is substituted or re-flowed.  Retained uncounted as the source-local context the proof needs: lines 1389--1393.  Nothing was omitted from any retained span, so the package carries no omission record.  No retained line contains a citation, so the wrapper carries no bibliography and no bibliography entry.  No retained line contains a figure environment, an image-inclusion command or drawing code, so no image is delivered and the package declares no asset dependency.  Editorial formatting only, in addition: an AMS \texttt{amsart} preamble that restates the article's own declarations and macros used by the retained lines, namely the environments \texttt{proposition}, \texttt{theorem} on separate counters, together with the standard packages those lines need.  The retained environments print in this document as Proposition 1, Theorem 1 in order of appearance, whereas the article prints the same items with the numbering of its own full document.  The complete native source of the article as distributed in the arXiv e-print for this exact version is bundled unchanged as \texttt{sources/167839/source0.tex}.
	\begin{proposition}\label{unitalcomm}
		The Grothendieck ring $\K^{\bi}(\lf)$ of a~Leibniz algebra $\lf$
		is commutative. Moreover, the class $[F_0^{s/a}]$ of the one-dimensional trivial $\lf$-bimodule $F_0^{s/a}$ is the unity of
		$\K^{\bi}(\lf)$.
	\end{proposition}

	\begin{theorem}\label{alt}
		Let $\lf$ be a~solvable Leibniz algebra over an algebraically closed field of characteristic zero whose canonical Lie algebra $\lf_{\lie}$ is finite dimensional. Then the Grothendieck ring $\K^{\bi}(\lf)$ is alternative.
	\end{theorem}

	\begin{proof}
		Since by Proposition~\ref{unitalcomm}, $\K^{\bi}(\lf)$ is commutative, it suffices to prove that $\K^{\bi}(\lf)$ is left alternative. By virtue of Lie's theorem, every finite-dimensional irreducible $\lf_{\lie}$-module is one-dimensional. Hence, we have that
		\[
			\irr(\lf_{\lie})=\{[F_\lambda]\mid\lambda\in\Lambda\}\,,
		\]
		where
		\[
			\Lambda:=\{\lambda\in\lf_{\lie}^*\mid\lambda([\lf_{\lie},\lf_{\lie}])=0\}\,,
		\]
		and therefore
		\[
			\irr^{\bi}(\lf)=\{[F_0^{s/a}]\}\cup\{[F_\lambda^s ]\mid\lambda\in\Lambda\setminus\{0\}\}
			\cup\{[F_\mu^a ]\mid\lambda\in\Lambda\setminus\{0\}\}\,.
		\]
		Let
		\[
			u=m_0 [F_0^{s/a}]+\sum_{\lambda\in\Lambda\setminus\{0\}}m_\lambda[F_\lambda^s ]
			+\sum_{\mu\in\Lambda\setminus\{0\}}m_\mu[F_\mu^a ]
		\]
		and
		\[
			v=n_0 [F_0^{s/a}]+\sum_{\nu\in\Lambda\setminus\{0\}}n_\nu[F_\nu^s ]
			+\sum_{\eta\in\Lambda\setminus\{0\}}n_\eta[F_\eta^a ]
		\]
		be arbitrary elements in $\K^{\bi}(\lf)$. Then we have that
		\begin{eqnarray*}
			u^2 & = & m_0^2 [F_0^{s/a}]+\sum_{\lambda\in\Lambda\setminus\{0\}}(2m_0 m_\lambda)
			[F_\lambda^s ]+\sum_{\mu\in\Lambda\setminus\{0\}}(2m_0 m_\mu)[F_\mu^a ]\\
			&& +\sum_{\lambda',\lambda''\in\Lambda\setminus\{0\}}(m_{\lambda'}m_{\lambda''})
			[F_{\lambda'+\lambda''}^s ]+\sum_{\mu',\mu''\in\Lambda\setminus\{0\}}
			(m_{\mu'}m_{\mu''})[F_{\mu'+\mu''}^a ]\,,
		\end{eqnarray*}
		and thus
		\begin{eqnarray*}
			u^2\cdot v & = & (m_0^2 n_0 )[F_0^{s/a}]+\sum_{\nu\in\Lambda\setminus\{0\}}(m_0^2 n_\nu)
			[F_\nu^s ]+\sum_{\eta\in\Lambda\setminus\{0\}}(m_0^2 n_\eta)[F_\eta^a ]\\
			&& +\sum_{\lambda\in\Lambda\setminus\{0\}}(2m_0 m_\lambda n_0 )[F_\lambda^s ]
			+\sum_{\lambda,\nu\in\Lambda\setminus\{0\}}(2m_0 m_\lambda n_\nu)[F_{\lambda+\nu}^s ]\\
			&& +\sum_{\mu\in\Lambda\setminus\{0\}}(2m_0 m_\mu n_0 )[F_\mu^a ]
			+\sum_{\mu,\eta\in\Lambda\setminus\{0\}}(2m_0 m_\mu n_\eta)[F_{\mu+\eta}^a ]\\
			&& +\sum_{\lambda',\lambda''\in\Lambda\setminus\{0\}}(m_{\lambda'}m_{\lambda''}n_0 )
			[F_{\lambda'+\lambda''}^s ]+\sum_{\lambda',\lambda'',\eta\in\Lambda\setminus\{0\}}
			(m_{\lambda'}m_{\lambda''}n_\nu)[F_{\lambda'+\lambda''+\nu}^s ]\\
			&& +\sum_{\mu',\mu''\in\Lambda\setminus\{0\}}(m_{\mu'}m_{\mu''}n_0 )[F_{\mu'+\mu''}^a ])
			+\sum_{\mu',\mu'',\eta\in\Lambda\setminus\{0\}}(m_{\mu'}m_{\mu''}n_\eta)[F_{\mu'+\mu''+\eta}^a ])\,.
		\end{eqnarray*}
		On the other hand, we have that
		\begin{eqnarray*}
			u\cdot v & = & (m_0 n_0 )[F_0^{s/a}]+\sum_{\nu\in\Lambda\setminus\{0\}}(m_0 n_\nu)[F_\nu^s ]
			+\sum_{\eta\in\Lambda\setminus\{0\}}(m_0 n_\eta)[F_\eta^a ]\\
			&& +\sum_{\lambda\in\Lambda\setminus\{0\}}(m_\lambda n_0 )[F_\lambda^s ]
			+\sum_{\lambda,\nu\in\Lambda\setminus\{0\}}(m_\lambda n_\nu)[F_{\lambda+\nu}^s ]\\
			&& +\sum_{\mu\in\Lambda\setminus\{0\}}(m_\mu n_0 )[F_\mu^a ]
			+\sum_{\mu,\eta\in\Lambda\setminus\{0\}}(m_\mu n_\eta)[F_{\mu+\eta}^a ]\,,
		\end{eqnarray*}
		and then
		\begin{eqnarray*}
			u\cdot(u\cdot v) & = & (m_0^2 n_0 )[F_0^{s/a}]+\sum_{\nu\in\Lambda\setminus\{0\}}(m_0^2 n_\nu)
			[F_\nu^s ]+\sum_{\eta\in\Lambda\setminus\{0\}}(m_0^2 n_\eta)[F_\eta^a ]\\
			&& +\underbrace{\sum_{\lambda\in\Lambda\setminus\{0\}}(m_0 m_\lambda n_0 )[F_\lambda^s ]}_{1}
			+\underbrace{\sum_{\lambda,\nu\in\Lambda\setminus\{0\}}(m_0 m_\lambda n_\nu)[F_{\lambda+\nu}^s ]}_{2}\\
			&& +\underbrace{\sum_{\mu\in\Lambda\setminus\{0\}}(m_0 m_\mu n_0 )[F_\mu^a ]}_{3}
			+\underbrace{\sum_{\mu,\eta\in\Lambda\setminus\{0\}}(m_0 m_\mu n_\eta)[F_{\mu+\eta}^a ]}_{4}\\
			&& +\underbrace{\sum_{\lambda\in\Lambda\setminus\{0\}}(m_\lambda m_0 n_0 )[F_\lambda^s ]}_{1}
			+\underbrace{\sum_{\lambda,\nu\in\Lambda\setminus\{0\}}(m_\lambda m_0 n_\nu)[F_{\lambda+\nu}^s ]}_{2}\\
			&& +\sum_{\lambda',\lambda''\in\Lambda\setminus\{0\}}(m_{\lambda'}m_{\lambda''}n_0 )
			[F_{\lambda'+\lambda''}^s ]
			+\sum_{\lambda',\lambda'',\nu\in\Lambda\setminus\{0\}}(m_{\lambda'}m_{\lambda''}n_\nu)
			[F_{\lambda'+\lambda''+\nu}^s ]\\
			&& +\underbrace{\sum_{\mu\in\Lambda\setminus\{0\}}(m_\mu m_0 n_0 )[F_\mu^a ]}_{3}
			+\underbrace{\sum_{\mu,\eta\in\Lambda\setminus\{0\}}(m_\mu m_0 n_\eta)[F_{\mu+\eta}^a ]}_{4}\\
			&& +\sum_{\mu',\mu''\in\Lambda\setminus\{0\}}(m_{\mu'}m_{\mu''}n_0 )[F_{\mu'+\mu''}^a ]
			+\sum_{\mu', \mu'', \eta\in\Lambda\setminus\{0\}}(m_{\mu'}m_{\mu''}n_\eta)[F_{\mu'+\mu''+\eta}^a ]\\
			& = & (m_0^2 n_0 )[F_0^{s/a}]+\sum_{\nu\in\Lambda\setminus\{0\}}(m_0^2 n_\nu)[F_\nu^s ]
			+\sum_{\eta\in\Lambda\setminus\{0\}}(m_0^2 n_\eta)[F_\eta^a ]\\
			&& +\sum_{\lambda\in\Lambda\setminus\{0\}}(2m_0 m_\lambda n_0 )[F_\lambda^s ]
			+\sum_{\lambda,\nu\in\Lambda\setminus\{0\}}(2m_0 m_\lambda n_\nu)[F_{\lambda+\nu}^s ]\\
			&& +\sum_{\mu\in\Lambda\setminus\{0\}}(2m_0 m_\mu n_0 )[F_\mu^a ]
			+\sum_{\mu,\eta\in\Lambda\setminus\{0\}}(2m_0 m_\mu n_\eta)[F_{\mu+\eta}^a ]\\
			&& +\sum_{\lambda',\lambda''\in\Lambda\setminus\{0\}}(m_{\lambda'}m_{\lambda''}n_0 )
			[F_{\lambda'+\lambda''}^s ]
			+\sum_{\lambda',\lambda'',\nu\in\Lambda\setminus\{0\}}(m_{\lambda'}m_{\lambda''}n_\nu)
			[F_{\lambda'+\lambda''+\nu}^s ]\\
			&& +\sum_{\mu',\mu''\in\Lambda\setminus\{0\}}(m_{\mu'}m_{\mu''}n_0 )[F_{\mu'+\mu''}^a ]
			+\sum_{\mu', \mu'', \eta\in\Lambda\setminus\{0\}}(m_{\mu'}m_{\mu''}n_\eta)[F_{\mu'+\mu''+\eta}^a ]\,,
		\end{eqnarray*}
		which completes the proof.
	\end{proof}

\end{document}
